% Phụ lục: ghi chú build. Nguồn: results/{apple-m5,LAP1184}/{campaign.log, machine.txt}, Makefile, mk/zz_pack.mk.

Trên \PCname{apple-m5}, đợt đo có 12\,152 lần chạy hoàn thành. Thời lượng 1 giờ 25 phút 45 giây gồm cả
2 phút hiệu chỉnh (643 lần chạy thử, không tính vào số trên). Lệnh đã dùng là
\texttt{python3 scripts/run\_bench.py --relax-yield --host apple-m5 --wait 30}, tương đương
\texttt{make bench} cộng lượt tùy chọn \texttt{make bench-relax}. \texttt{make bench-dry} in lưới và thời
gian ước tính mà không đo.
% nguồn: campaign.log (START argv=..., END duration=1:25:45); results/apple-m5/campaign_stdout.log ("calibrating 643
% iteration counts", "calibration done in 2:04")

Hai thí nghiệm phụ trên \PCname{apple-m5} (khoảng trống khi trao lock và đo lại lúc tải thấp) dùng
\texttt{experiments/handoff\_gap.cpp} và \texttt{experiments/run\_extra.py}. Chương trình
\texttt{handoff\_gap} chỉ cần các header trong \texttt{include/}, build bằng \texttt{mkdir -p
build/experiments \&\& c++ -Iinclude -std=c++20 -O2 -pthread experiments/handoff\_gap.cpp -o
build/experiments/handoff\_gap}. Lượt đo lại dùng \texttt{build/bench\_lock}, nên cần chạy
\texttt{make} trước. Lệnh \texttt{python3 experiments/run\_extra.py apple-m5} ghi
\texttt{handoff.csv}, \texttt{lowload.csv} và \texttt{run.log} vào \texttt{results/apple-m5/extra/}.
Số trong báo cáo tính bằng \texttt{scripts/report\_numbers/extra\_analyze.py}.
% nguồn: experiments/handoff_gap.cpp (include: bench_common.hpp, lock.h, platform.hpp, workload.hpp;
% omp.h chỉ khi _OPENMP); experiments/run_extra.py (docstring: Part A, Part B, Output)

Nền tảng gồm macOS 27.0.1, Apple clang 21.0.0, Python 3.14.8 (chạy đo chỉ cần thư viện chuẩn).
Apple clang không nhận \texttt{-fopenmp} nên không có baseline OpenMP. \texttt{matplotlib} và
\texttt{numpy} cài trong một môi trường ảo (\texttt{.venv}).
% nguồn: results/apple-m5/machine.txt

Trên \PCname{LAP1184}, nhóm chạy toàn bộ quy trình bằng \texttt{make full PYTHON=.venv/bin/python}
trong Ubuntu 20.04 trên WSL2. Đợt đo (\texttt{run\_bench.py --wait 10}) dài 1 giờ 20 phút 36 giây
(14\,322 lần chạy hoàn thành). Ubuntu 20.04 cần thêm g++ 11 từ PPA \texttt{ubuntu-toolchain-r/test} (g++ 9 mặc
định không có \texttt{<barrier>}). Nó cũng cần CMake 3.30 từ bản cài sẵn của Kitware (bản 3.16 của Ubuntu dưới mức
3.20 mà \texttt{CMakeLists.txt} cần), \texttt{python3-venv}, XeLaTeX và \texttt{latexmk}. Python 3.8.10
đủ cho các script.
\begin{itemize}[nosep]
  \item Git trên Windows mặc định đổi kết thúc dòng thành CRLF khi checkout, và \texttt{sh} trong WSL
    không chạy được \texttt{scripts/machine\_info.sh} khi đó. Runner bắt lỗi và ghi một
    \texttt{machine.txt} rỗng. Kho mã có \texttt{.gitattributes} giữ LF cho các file \texttt{.sh}.
    \texttt{machine.txt} của \PCname{LAP1184} được sinh lại bằng đúng lệnh của runner, 3 phút sau khi đợt
    đo kết thúc. Vì vậy dòng nguồn điện và load average trong đó là của lúc sinh lại.
  \item ThreadSanitizer báo data race giả trên \texttt{omp\_lock\_t}, vì libgomp không được instrument.
    Bản TSan của \texttt{test\_locks} và \texttt{test\_counter} bỏ qua lock OpenMP và in ra việc bỏ
    qua, như \texttt{test\_barrier} làm với \texttt{omp\_barrier}. Tham số \texttt{tsan\_uninstrumented}
    vẫn chạy nó để thấy cảnh báo.
  \item Không chạy \texttt{make bench-relax}, vì lựa chọn \texttt{isb}/\texttt{yield} chỉ có trên AArch64.
\end{itemize}
% nguồn: results/LAP1184/campaign.log (START argv=--wait 10; END duration=1:20:36, 01:41:24);
% results/LAP1184/machine.txt (date: 01:44:04); docs/part1.md mục 7.2.
Trên \PCname{C4M7Y670NT}, đợt đo (\texttt{run\_bench.py --wait 10}) dài 1 giờ 53 phút 10 giây (13\,244 lần
chạy hoàn thành), với macOS 26.7.1, Apple clang 21.0.0 và Python 3.9.6. Cách cài giống \PCname{apple-m5}.
% nguồn: results/C4M7Y670NT/campaign.log (START argv=--wait 10; END duration=1:53:10 ok=13244);
% results/C4M7Y670NT/machine.txt (python: Python 3.9.6).
